\documentclass[aspectratio=169,11pt,english]{beamer}
\input{../../../_preambulo_beamer.tex}
\input{../../../_comandos_beamer.tex}

\newcommand{\ruta}{Class route · 40 minutes}
\title[L05 · \ruta]{Set Theory and Sample Spaces}
% Code fragments support the explanation; source files remain the authoritative
% place to read every docstring and comment at full size.
\lstset{basicstyle=\ttfamily\tiny,columns=fullflexible,keepspaces=true}
\newcommand{\archivo}[1]{{\footnotesize\texttt{\detokenize{#1}}\par}\smallskip}
\newcommand{\objetivo}[1]{{\small\color{TecAzul}#1\par}\medskip}
\newcommand{\question}[1]{\par\smallskip\textbf{Your task.} #1\par}
\newcommand{\checkit}[1]{\par\smallskip\textbf{Check.} #1\par}
\newcommand{\code}[3]{%
\ifnum#2>50\lstset{gobble=4}\else\lstset{gobble=0}\fi
\ifcase#1\relax
\or\lstinputlisting[firstline=#2,lastline=#3]{../src/set_theory_01_sample_space.py}%
\or\lstinputlisting[firstline=#2,lastline=#3]{../src/set_theory_02_set_operations.py}%
\or\lstinputlisting[firstline=#2,lastline=#3]{../src/set_theory_03_disjoint_events.py}%
\fi}

\begin{document}

\begin{frame}{Set Theory and Sample Spaces}
  \framesubtitle{Class route · 40 minutes}
  \objetivo{L05 · A student guide to mathematics and programming}
  How can we identify night dispatches, delayed dispatches, and dispatches that
  satisfy both conditions?
  \begin{block}{What you will build}
    A model of possible outcomes and a Python program that answers questions
    with sets, membership, and reproducible observations.
  \end{block}
  Synthetic dispatch case: context, definitions, implementation, and checks.
  \par\medskip
  {\footnotesize MA1001B Faculty \textbar{} Course lessons \textbar{} Tecnol\'ogico de Monterrey}
\end{frame}

\begin{frame}{By the end, you can demonstrate that you can...}
  \framesubtitle{Class route · 40 minutes}
  \begin{enumerate}
    \item Define an experiment, represent an outcome, and construct its sample space $S$.
    \item Distinguish possible outcomes, observed records, cardinalities, and frequencies.
    \item Represent events as subsets and calculate union, intersection, and complement.
    \item Justify disjointness and exhaustiveness as separate properties.
    \item Implement the operations in Python and interpret their results.
  \end{enumerate}
  \checkit{Every objective has an activity and a check; the final question combines all five.}
\end{frame}

\begin{frame}{Your learning route: 40 minutes}
  \framesubtitle{Class route · 40 minutes}
  \begin{tabular}{p{.34\linewidth}rp{.45\linewidth}}
    \toprule Block & Min & Product \\ \midrule
    Case and objectives & 5 & Interpret one dispatch.\\
    Sample space & 10 & Construct eight possibilities.\\
    Events and operations & 15 & Solve and program set questions.\\
    Disjointness and exhaustiveness & 11 & Justify two independent properties.\\
    Integration & 4 & Solve a new question.\\ \bottomrule
  \end{tabular}
  \par\medskip
  Follow intuition $\to$ formal definition $\to$ worked example $\to$ implementation $\to$ interpretation.
  The simulation details and additional proofs appear in the student consultation section.
\end{frame}

\begin{frame}{The case: classify one dispatch}
  \framesubtitle{Class route · 40 minutes}
  \objetivo{Objective 1 · Context before symbols}
  A \textbf{dispatch} is one individual shipment. We record:
  \begin{description}
    \item[Shift] Day or Night: when the dispatch is processed.
    \item[Service] Standard or Express: the assigned service level.
    \item[Status] On time or Delayed: whether the delivery met its target.
  \end{description}
  These categories let us ask about time, service, and delivery performance.
  \begin{alertblock}{Synthetic teaching model}
    The categories, 480 records, and generating probabilities are invented.
    They are not observations or estimates from a real company.
  \end{alertblock}
\end{frame}

\begin{frame}{Experiment and outcome: observe one repetition}
  \framesubtitle{Class route · 40 minutes}
  \textbf{Experiment:} observe one dispatch and record its three characteristics;
  the combination is not known in advance.
  \[
    \omega=(\text{shift},\text{service},\text{status})
  \]
  The symbol $\omega$ represents one possible outcome. For example,
  \[
    \omega=(\text{Night},\text{Express},\text{Delayed}).
  \]
  This means a night dispatch with Express service that missed its delivery target.
  \question{Why is writing only ``Delayed'' insufficient?}
  \checkit{Shift and service are missing; position is part of the outcome contract.}
\end{frame}

\begin{frame}{A row, an identifier, and a result}
  \framesubtitle{Class route · 40 minutes}
  \begin{tabular}{rlll}
    \toprule ID & Shift & Service & Status \\ \midrule
    101 & Day & Standard & On time\\
    102 & Day & Standard & On time\\
    103 & Night & Express & Delayed\\ \bottomrule
  \end{tabular}
  \par\medskip
  These three hypothetical shipments contain \textbf{two distinct outcomes}.
  The ID distinguishes records; it is not part of $\omega$.
  \[
    3\ \text{records}\qquad 2\ \text{observed outcomes}\qquad |S|=8.
  \]
  \question{Does the second row increase $|S|$?}
  \checkit{No. It repeats one outcome and increases its frequency, not the possibilities.}
\end{frame}

\begin{frame}{Build the sample space}
  \framesubtitle{Class route · 40 minutes}
  Define $T=\{\text{Day},\text{Night}\}$, $V=\{\text{Standard},\text{Express}\}$,
  and $E=\{\text{On time},\text{Delayed}\}$.
  \[
    S=T\times V\times E=\{(t,v,e):t\in T,\ v\in V,\ e\in E\}.
  \]
  The \textbf{Cartesian product} forms all tuples by choosing one element from
  every set. The symbol $\in$ means ``belongs to.''
  \[
    |S|=|T|\,|V|\,|E|=2\cdot2\cdot2=8.
  \]
  $|S|$ is the \textbf{cardinality}: the number of distinct possible outcomes.
  Possibility does not imply equal probability.
\end{frame}

\begin{frame}{The eight possibilities, one per row}
  \framesubtitle{Class route · 40 minutes}
  \centering
  \begin{tabular}{clll}
    \toprule Label & Shift & Service & Status \\ \midrule
    $\omega_1$ & Day & Standard & On time\\
    $\omega_2$ & Day & Standard & Delayed\\
    $\omega_3$ & Day & Express & On time\\
    $\omega_4$ & Day & Express & Delayed\\
    $\omega_5$ & Night & Standard & On time\\
    $\omega_6$ & Night & Standard & Delayed\\
    $\omega_7$ & Night & Express & On time\\
    $\omega_8$ & Night & Express & Delayed\\ \bottomrule
  \end{tabular}
  \par\medskip\raggedright
  $\omega_i$ abbreviates a \textbf{type of outcome}; it does not identify one shipment.
  \question{List the four Day possibilities and then the four Night possibilities.}
\end{frame}

\begin{frame}{Increment 1 · Enumerate without omissions}
  \framesubtitle{Class route · 40 minutes}
  \archivo{set_theory_01_sample_space.py}
  \textbf{Function:} \texttt{build\_sample\_space()}.
  \textbf{Available inputs:} the three category tuples.
  \code{1}{57}{61}
  \question{Combine one category from each collection with \texttt{product};
  preserve the tuple order so it can be checked against the table.}
  \checkit{Eight distinct tuples. First is (Day, Standard, On time); last is
  (Night, Express, Delayed).}
\end{frame}

\begin{frame}{Increment 1 · Read the product contract}
  \framesubtitle{Class route · 40 minutes}
  \archivo{set_theory_01_sample_space.py}
  \code{1}{62}{71}
  The last argument changes fastest, so the order is predictable. Materializing
  the iterator lets later tables and figures reuse the same sequence.
  \checkit{The return value has eight inner tuples, each with three strings.}
\end{frame}

\begin{frame}{Increment 1 · Materialize and verify}
  \framesubtitle{Class route · 40 minutes}
  \archivo{set_theory_01_sample_space.py}
  \code{1}{65}{71}
  \texttt{product} replaces three manually nested loops. The outer \texttt{tuple}
  materializes the iterator and keeps one stable order; each inner tuple is one result.
  \[
    |S|=8\qquad |\{\omega:\omega\text{ is listed twice}\}|=0.
  \]
  \checkit{Test both \texttt{len(sample\_space) == 8} and
  \texttt{len(set(sample\_space)) == 8}. Eight positions alone would not rule out duplicates.}
\end{frame}

\begin{frame}{From eight possibilities to 480 observations}
  \framesubtitle{Class route · 40 minutes}
  \archivo{set_theory_01_sample_space.py}
  \texttt{simulate\_dispatches()} builds a DataFrame in memory with seed 42.
  It does not download data or write a CSV.
  \begin{itemize}
    \item Shift and service use unequal generating probabilities.
    \item The delay probability depends on shift and service.
    \item The larger delay increment for Express is a synthetic assumption,
          not a claim about Express in reality.
  \end{itemize}
  This run contains all eight outcome types. The most frequent type appears 196 times.
  \checkit{Possible, observed, and frequent are different statements.}
\end{frame}

\begin{frame}{Increment 1 · Generate aligned records}
  \framesubtitle{Class route · 40 minutes}
  \archivo{set_theory_01_sample_space.py}
  \textbf{Function:} \texttt{simulate\_dispatches(seed)}.
  \code{1}{75}{82}
  \question{First draw one shift and one service for each of the 480 records.
  Read the probability tuples in category order.}
  \checkit{Both arrays have length 480; probabilities describe the model and do
  not impose exact category quotas.}
\end{frame}

\begin{frame}{Increment 1 · One probability per record}
  \framesubtitle{Class route · 40 minutes}
  \archivo{set_theory_01_sample_space.py}
  \code{1}{108}{119}
  A Boolean comparison contributes 1 when its condition is true and 0 otherwise.
  Thus the base probability and the two increments produce 0.10, 0.18, 0.22,
  or 0.30 for a row.
  \checkit{Every delay probability lies in $[0,1]$; then one uniform draw chooses
  the status for each row.}
\end{frame}

\begin{frame}{Increment 1 · Finalize each status}
  \framesubtitle{Class route · 40 minutes}
  \archivo{set_theory_01_sample_space.py}
  \code{1}{120}{130}
  The comparison now produces one status per row.
\end{frame}

\begin{frame}{Increment 1 · Assemble the observed table}
  \framesubtitle{Class route · 40 minutes}
  \archivo{set_theory_01_sample_space.py}
  \code{1}{131}{138}
  \texttt{dispatch\_id} labels a row. \texttt{zip} joins the aligned shift,
  service, and status values into one outcome tuple.
\end{frame}

\begin{frame}{Increment 1 · Assemble the observed table}
  \framesubtitle{Class route · 40 minutes}
  \archivo{set_theory_01_sample_space.py}
  \code{1}{139}{145}
  The tuple column preserves the three-component result while the ID keeps
  records distinguishable.
  \checkit{The output has 480 rows and columns \texttt{dispatch\_id},
  \texttt{shift}, \texttt{service}, \texttt{status}, and \texttt{outcome}.}
\end{frame}

\begin{frame}{Read frequencies without losing possibilities}
  \framesubtitle{Class route · 40 minutes}
  \centering
  \includegraphics[height=.64\textheight,width=.96\linewidth,keepaspectratio]{../figures/set_theory_01_sample_space.png}
  \par\raggedright{\small Each bar counts synthetic records of one possible tuple;
  all bars sum to 480. The bar of 196 does not add possibilities to $S$.}
\end{frame}

\begin{frame}{Events select outcomes from $S$}
  \framesubtitle{Class route · 40 minutes}
  \objetivo{Objective 3 · Question: which types are Night or Delayed?}
  An \textbf{event} is a subset of $S$: it selects outcomes satisfying a condition.
  \[
    A=\{\omega\in S:\text{shift = Night}\}=\{\omega_5,\omega_6,\omega_7,\omega_8\}
  \]
  \[
    B=\{\omega\in S:\text{status = Delayed}\}=\{\omega_2,\omega_4,\omega_6,\omega_8\}.
  \]
  $A\subseteq S$ means every element of $A$ belongs to $S$.
  \checkit{$|A|=|B|=4$ possible types; this is not yet a count of observed rows.}
\end{frame}

\begin{frame}{Increment 2 · Select by condition}
  \framesubtitle{Class route · 40 minutes}
  \archivo{set_theory_02_set_operations.py}
  \textbf{Function:} \texttt{define\_events(sample\_space)}.
  \question{Build two sets: one for Night and one for Delayed. Filter the eight
  possible tuples, not the 480 records.}
  \begin{itemize}
    \item Position 0 is shift; position 1 is service; position 2 is status.
    \item Keep the complete tuple when its component satisfies the condition.
  \end{itemize}
  \code{2}{92}{96}
  \checkit{The set comprehension begins with the first event condition.}
\end{frame}

\begin{frame}{Increment 2 · Define the second event}
  \framesubtitle{Class route · 40 minutes}
  \archivo{set_theory_02_set_operations.py}
  \code{2}{97}{104}
  \checkit{Each set has four elements; $\omega_6$ and $\omega_8$ belong to both.}
\end{frame}

\begin{frame}{Increment 2 · Keep the complete tuple}
  \framesubtitle{Class route · 40 minutes}
  \archivo{set_theory_02_set_operations.py}
  \code{2}{100}{104}
  The condition examines one component but returns the whole outcome. The event
  remains a subset of $S$, not a collection of isolated labels.
\end{frame}

\begin{frame}{Intersection: both conditions}
  \framesubtitle{Class route · 40 minutes}
  \textbf{Question:} which possible types are Night \textbf{and} Delayed?
  \[
    A\cap B=\{\omega\in S:\omega\in A\ \text{and}\ \omega\in B\}.
  \]
  From $A$, $\omega_5$ and $\omega_7$ are On time, so discard them:
  \[
    A\cap B=\{\omega_6,\omega_8\},\qquad |A\cap B|=2.
  \]
  They are (Night, Standard, Delayed) and (Night, Express, Delayed).
  \checkit{The result is two possible types, not necessarily two observed dispatches.}
\end{frame}

\begin{frame}{Union: at least one condition}
  \framesubtitle{Class route · 40 minutes}
  \textbf{Question:} which possible types are Night \textbf{or} Delayed?
  Here ``or'' is inclusive: both conditions are accepted.
  \[
    A\cup B=\{\omega\in S:\omega\in A\ \text{or}\ \omega\in B\}.
  \]
  \[
    A\cup B=\{\omega_2,\omega_4,\omega_5,\omega_6,\omega_7,\omega_8\}.
  \]
  The shared outcomes $\omega_6$ and $\omega_8$ enter once each.
  \checkit{Six types qualify; only Day and On time types $\omega_1,\omega_3$ are absent.}
\end{frame}

\begin{frame}{Why $4+4-2=6$: inclusion--exclusion}
  \framesubtitle{Class route · 40 minutes}
  \begin{tabular}{lll}
    \toprule Region & Outcomes & Count \\ \midrule
    Only $A$ & $\omega_5,\omega_7$ & 2\\
    $A\cap B$ & $\omega_6,\omega_8$ & 2\\
    Only $B$ & $\omega_2,\omega_4$ & 2\\ \bottomrule
  \end{tabular}
  \par\medskip
  Adding $4+4$ counts $\omega_6$ twice and $\omega_8$ twice. Subtract one copy
  of the shared region:
  \[
    |A\cup B|=|A|+|B|-|A\cap B|=4+4-2=6.
  \]
  \checkit{The formula counts every member of the union exactly once.}
\end{frame}

\begin{frame}{Complement: what is missing inside the universe}
  \framesubtitle{Class route · 40 minutes}
  \textbf{Question:} which possible types are not Night?
  The universe remains $S$ and $A$ contains the Night outcomes.
  \[
    A^c=S\setminus A=\{\omega\in S:\omega\notin A\}
      =\{\omega_1,\omega_2,\omega_3,\omega_4\}.
  \]
  These are the four Day types. The superscript $c$ means complement, not a power.
  \begin{alertblock}{The universe matters}
    We do not add a Dawn or Midnight shift: it is not in this $S$.
    A complement is always relative to a stated universe.
  \end{alertblock}
  \checkit{$|A^c|=8-4=4$ possible types.}
\end{frame}

\begin{frame}{Increment 2 · Program the operations}
  \framesubtitle{Class route · 40 minutes}
  \archivo{set_theory_02_set_operations.py}
  \textbf{Function:} \texttt{summarize\_set\_operations()}.
  \question{Validate that $A$ and $B$ are subsets of $S$, then return the five named sets.}
  \begin{tabular}{lll}
    \toprule Operation & Python & Expected size \\ \midrule
    $A$ and $B$ & stored directly & 4 and 4\\
    Intersection & \texttt{\&} & 2\\
    Union & \texttt{|} & 6\\
    Complement & \texttt{set(S) - A} & 4\\ \bottomrule
  \end{tabular}
  The table states the intended contract before the implementation.
\end{frame}

\begin{frame}{Increment 2 · Validate the universe}
  \framesubtitle{Class route · 40 minutes}
  \archivo{set_theory_02_set_operations.py}
  \code{2}{107}{112}
  \question{Start with the universe and reject an event containing a tuple that
  cannot occur under the stated model.}
\end{frame}

\begin{frame}{Increment 2 · Apply the validated operations}
  \framesubtitle{Class route · 40 minutes}
  \archivo{set_theory_02_set_operations.py}
  \code{2}{113}{118}
  Validation makes the relative complement unambiguous: it is computed against
  the universe received by the function.
\end{frame}

\begin{frame}{Increment 2 · Apply the three set operators}
  \framesubtitle{Class route · 40 minutes}
  \archivo{set_theory_02_set_operations.py}
  \code{2}{119}{134}
  \texttt{\&} means intersection, \texttt{|} means union, and subtraction from
  the universe means relative complement.
\end{frame}

\begin{frame}{Increment 2 · Name the returned operations}
  \framesubtitle{Class route · 40 minutes}
  \archivo{set_theory_02_set_operations.py}
  \code{2}{127}{134}
  The dictionary names answers so later code can inspect them without repeating
  the same set expression. It still counts possible types, not rows.
  \checkit{The five lengths are 4, 4, 2, 6, and 4.}
\end{frame}

\begin{frame}{Now count observed records}
  \framesubtitle{Class route · 40 minutes}
  A set answers ``which types are possible?'' A Boolean mask answers ``which
  rows in this DataFrame satisfy the condition?''
  \[
    n(A)=172,\quad n(B)=86,\quad n(A\cap B)=45,\quad n(A\cup B)=213.
  \]
  These are counts of rows, not cardinalities of event sets. The same logic gives
  \[
    n(A\cup B)=172+86-45=213,\qquad n(A^c)=480-172=308.
  \]
  \question{Use pandas \texttt{\&} and \texttt{|} between masks; \texttt{and} and
  \texttt{or} do not perform elementwise Boolean operations on a Series.}
\end{frame}

\begin{frame}{Cardinality and frequency use different units}
  \framesubtitle{Class route · 40 minutes}
  \begin{tabular}{lrr}
    \toprule Event & $|G|$: possible types & $n(G)$: observed dispatches \\ \midrule
    $A$ & 4 & 172\\
    $B$ & 4 & 86\\
    $A\cap B$ & 2 & 45\\
    $A\cup B$ & 6 & 213\\
    $A^c$ & 4 & 308\\ \bottomrule
  \end{tabular}
  \par\medskip
  A type may repeat in many rows. That is why $|A\cap B|=2$ and
  $n(A\cap B)=45$ can both be correct.
  \checkit{State the unit every time: possible types or observed dispatches.}
\end{frame}

\begin{frame}{A membership matrix, not a frequency chart}
  \framesubtitle{Class route · 40 minutes}
  \centering
  \includegraphics[height=.64\textheight,width=.96\linewidth,keepaspectratio]{../figures/set_theory_02_set_operations.png}
  \par\raggedright{\small Each cell is 1 if a possible outcome belongs to a set and 0 otherwise.
  The figure does not represent the 480 observed rows.}
\end{frame}

\begin{frame}{Two questions you must not confuse}
  \framesubtitle{Class route · 40 minutes}
  \objetivo{Objective 4 · Disjointness and exhaustiveness}
  \textbf{Can both events occur together?} They are \textbf{disjoint} when
  \[
    A\cap B=\varnothing.
  \]
  \textbf{Do they cover every possibility?} They are \textbf{exhaustive} with
  respect to $S$ when
  \[
    A\cup B=S.
  \]
  Exhaustive events may overlap. The tests answer different questions.
\end{frame}



\begin{frame}{Increment 3 · Test properties independently}
  \framesubtitle{Class route · 40 minutes}
  \archivo{set_theory_03_disjoint_events.py}
  \textbf{Functions:} \texttt{define\_service\_events()} and
  \texttt{compare\_event\_pairs()}.
  \question{Build C/D, then compare A/B and C/D with two independent tests.}
  \code{3}{93}{100}
  \checkit{C and D each filter the service component (position 1) and form a
  partition of S.}
\end{frame}

\begin{frame}{Increment 3 · Complete the service events}
  \framesubtitle{Class route · 40 minutes}
  \archivo{set_theory_03_disjoint_events.py}
  \code{3}{101}{106}
  Both event definitions use the same universe and differ only in the service
  label. This is why the pair partitions S.
\end{frame}

\begin{frame}{Increment 3 · Compare each pair}
  \framesubtitle{Class route · 40 minutes}
  \archivo{set_theory_03_disjoint_events.py}
  \code{3}{108}{115}
  \question{For each named pair, validate membership in S and calculate the
  intersection and union exactly once.}
\end{frame}

\begin{frame}{Increment 3 · Record the comparison}
  \framesubtitle{Class route · 40 minutes}
  \archivo{set_theory_03_disjoint_events.py}
  \code{3}{116}{123}
  The comparison begins with an empty row list and one loop over the named pairs.
\end{frame}

\begin{frame}{Increment 3 · Record the comparison}
  \framesubtitle{Class route · 40 minutes}
  \archivo{set_theory_03_disjoint_events.py}
  \code{3}{124}{138}
  Reusing the calculated intersection and union keeps the two property checks
  DRY and makes their different meanings visible.
  \checkit{One row per pair with separate property columns.}
\end{frame}

\begin{frame}{Increment 3 · Translate definitions into tests}
  \framesubtitle{Class route · 40 minutes}
  \archivo{set_theory_03_disjoint_events.py}
  \code{3}{139}{151}
  \texttt{not intersection} tests emptiness; \texttt{union == universe} tests
  coverage.
\end{frame}

\begin{frame}{Increment 3 · Keep the tests independent}
  \framesubtitle{Class route · 40 minutes}
  \archivo{set_theory_03_disjoint_events.py}
  \code{3}{152}{160}
  Equal sizes alone would not prove set equality, so the code compares the union
  with the received universe.
  \checkit{A/B: intersection 2, union 6, false/false. C/D: intersection 0,
  union 8, true/true.}
\end{frame}

\begin{frame}{Visual comparison: overlap versus coverage}
  \framesubtitle{Class route · 40 minutes}
  \centering
  \includegraphics[height=.64\textheight,width=.96\linewidth,keepaspectratio]{../figures/set_theory_03_disjoint_events.png}
  \par\raggedright{\small The pink bars answer the overlap question, the navy bars answer
  the coverage question, and the dashed line is $|S|=8$.}
\end{frame}

\begin{frame}{Integration · A new question}
  \framesubtitle{Class route · 40 minutes}
  \objetivo{Objective 5 · Combine condition, set operation, code, and interpretation}
  \textbf{Context:} identify Express dispatches that arrived Delayed.
  \textbf{Question:} what possible types satisfy both conditions, how many records
  were observed, and how does this event compare with $A^c$ (Day)?
  Define
  \[
    F=\{r\in S:r_1=\text{Express},\ r_2=\text{Delayed}\}.
  \]
  \question{Before reading the answer, write the two members of $F$, estimate its
  observed count, and predict $F\cap A^c$ and $F\cup A^c$.}
\end{frame}



\begin{frame}{Check the five learning objectives}
  \framesubtitle{Class route · 40 minutes}
  \begin{enumerate}
    \item Can you explain the experiment and enumerate $S$ without the table?
    \item Can you explain why $|S|=8$ even though there are 480 rows?
    \item Can you justify the members of an intersection, union, and complement?
    \item Can you test disjointness and exhaustiveness independently?
    \item Can you predict, run, and interpret the Python results?
  \end{enumerate}
  \checkit{The next lesson counts larger spaces without listing every outcome.}
\end{frame}

\begin{frame}{Reading and practice}
  \framesubtitle{Class route · 40 minutes}
  Holmes, Illowsky, and Dean, \emph{Introductory Business Statistics 2e}, OpenStax:
  \begin{itemize}
    \item \href{https://openstax.org/books/introductory-business-statistics-2e/pages/3-1-terminology}{Section 3.1: Terminology}. Review experiments, outcomes, sample spaces, events, unions, intersections, and complements.
    \item \href{https://openstax.org/books/introductory-business-statistics-2e/pages/3-2-independent-and-mutually-exclusive-events}{Section 3.2: Independent and Mutually Exclusive Events}. Contrast mutual exclusivity with independence.
  \end{itemize}
  The conceptual foundation is CC BY-NC-SA 4.0. The dispatch case, code,
  figures, and numerical results are original course material using synthetic data.
\end{frame}

\appendix
\renewcommand{\ruta}{Student consultation}

\begin{frame}{Student consultation · Choose what to revisit}
  \framesubtitle{Student consultation}
  \begin{itemize}
    \item A worked solution for $F$ and counterexamples.
    \item The general inclusion--exclusion argument.
    \item Why disjointness is not independence.
    \item Simulation assumptions, seed, and data structures.
    \item Declarative graphics, resource ownership, and validation.
    \item Additional practice with reasoned solutions.
  \end{itemize}
  These pages are part of the student guide; they are not instructor notes.
\end{frame}

\begin{frame}{Inclusion--exclusion: the general reason}
  \framesubtitle{Student consultation}
  Split the union into three non-overlapping regions:
  \[
    A\setminus B,\qquad A\cap B,\qquad B\setminus A.
  \]
  If their cardinalities are $x,y,z$, then
  \[
    |A|=x+y,\quad |B|=y+z,\quad |A\cup B|=x+y+z.
  \]
  The sum $|A|+|B|=x+2y+z$ counts the shared region twice. Subtract one copy:
  \[
    \boxed{|A\cup B|=|A|+|B|-|A\cap B|}.
  \]
  In this case $x=y=z=2$.
\end{frame}

\begin{frame}{The properties combine in four ways}
  \framesubtitle{Student consultation}
  \begin{tabular}{llll}
    \toprule Pair & $|\cap|$ & $|\cup|$ & Conclusion \\ \midrule
    $C,D$ & 0 & 8 & Disjoint and exhaustive\\
    $A,B$ & 2 & 6 & Neither property\\
    $A, A^c\cap B$ & 0 & 6 & Disjoint, not exhaustive\\
    $S,A$ & 4 & 8 & Exhaustive, not disjoint\\ \bottomrule
  \end{tabular}
  \par\medskip
  A property does not imply the other. For example, $S$ and $A$ cover $S$ but overlap.
\end{frame}

\begin{frame}{Disjointness does not mean independence}
  \framesubtitle{Student consultation}
  Independence is a probability statement:
  \[
    P(A\cap C)=P(A)P(C).
  \]
  In the generating model, shift and service are independently sampled:
  $P(A)=0.35$ and $P(C)=0.70$, so $P(A\cap C)=0.245$ even though $A$ and $C$
  share possible outcomes.
  Two disjoint events with positive probability are not independent: their
  intersection has probability zero but the product is positive.
  \checkit{Set relationships and probability relationships are different claims.}
\end{frame}

\begin{frame}{The universe determines a complement}
  \framesubtitle{Student consultation}
  Suppose Express does not operate at Night. This is a new model, not a plotting edit:
  \[
    S'=S\setminus\{\omega_7,\omega_8\},\qquad |S'|=6.
  \]
  Events must be rebuilt as subsets of $S'$, and the simulator must respect the new
  restriction. The complement of Night in $S'$ is then the Day outcomes only.
  \checkit{Reusable functions should use the universe received, not a hard-coded 8.}
\end{frame}

\begin{frame}{Simulation · probabilities are assumptions}
  \framesubtitle{Student consultation}
  \begin{tabular}{llr}
    \toprule Shift & Service & Delay probability \\ \midrule
    Day & Standard & 0.10\\
    Night & Standard & 0.18\\
    Day & Express & 0.22\\
    Night & Express & 0.30\\ \bottomrule
  \end{tabular}
  \par\medskip
  The 0.30 is $0.10+0.08+0.12$. It is a transparent teaching assumption.
  It does not estimate the reliability of Express in a real operation.
  \checkit{The four model rates are 0.10, 0.18, 0.22, and 0.30.}
\end{frame}

\begin{frame}{Simulation · apply the two increments}
  \framesubtitle{Student consultation}
  \archivo{set_theory_01_sample_space.py}
  \code{1}{118}{123}
  The comparisons are arrays aligned with the 480 dispatches. The validation
  rejects an impossible probability rather than silently clipping it.
\end{frame}

\begin{frame}{Simulation · apply the two increments}
  \framesubtitle{Student consultation}
  \archivo{set_theory_01_sample_space.py}
  \code{1}{124}{129}
  The Night and Express indicators contribute 0 or their named increments to
  the base rate. This is vectorized arithmetic, not a row-by-row loop.
\end{frame}

\begin{frame}{Simulation · Convert probability to status}
  \framesubtitle{Student consultation}
  \archivo{set_theory_01_sample_space.py}
  \code{1}{129}{136}
  For a row with probability $p$, \texttt{np.where} implements
  \textit{Delayed if $U<p$; On time otherwise}.
\end{frame}

\begin{frame}{Simulation · assemble the tuple column}
  \framesubtitle{Student consultation}
  \archivo{set_theory_01_sample_space.py}
  \code{1}{137}{145}
  \texttt{zip} joins values with the same row position. Repeated tuples remain
  repeated observations even though the tuple itself is immutable.
\end{frame}

\begin{frame}{Data structures with a purpose}
  \framesubtitle{Student consultation}
  \begin{description}
    \item[Inner tuple] An ordered result: shift, service, status. Immutable and hashable.
    \item[Outer tuple] The eight possibilities in a stable display order.
    \item[Set] An event without duplicates, with membership and set operations.
    \item[DataFrame] The 480 records, including repeated outcomes and IDs.
  \end{description}
  Turning observed outcomes into a set answers ``which types appeared?'' but
  removes the repetitions needed to calculate frequencies.
\end{frame}

\begin{frame}{Graphics · delegate presentation, retain reasoning}
  \framesubtitle{Student consultation}
  \archivo{set_theory_01_sample_space.py}
  \textbf{Function:} \texttt{build\_frequency\_figure()}.
  \code{1}{174}{181}
  The function first builds an explicit table in S order. This is the academic
  calculation: each row is one possible outcome and one observed count.
\end{frame}

\begin{frame}{Graphics · let Seaborn handle geometry}
  \framesubtitle{Student consultation}
  \archivo{set_theory_01_sample_space.py}
  \code{1}{182}{188}
  The frequency list uses \texttt{get(item, 0)} so a possible but absent outcome
  remains visible.
\end{frame}

\begin{frame}{Graphics · let Seaborn handle geometry}
  \framesubtitle{Student consultation}
  \archivo{set_theory_01_sample_space.py}
  \code{1}{189}{198}
  \texttt{barplot} places semantic categories and bars; the preceding table still
  computes the count for each member of $S$.
\end{frame}

\begin{frame}{Graphics · label exact counts}
  \framesubtitle{Student consultation}
  \archivo{set_theory_01_sample_space.py}
  \code{1}{199}{208}
  \texttt{errorbar=None} avoids inferential intervals that this example did not
  estimate; \texttt{bar\_label} resolves bar coordinates.
\end{frame}

\begin{frame}{Graphics · label exact counts}
  \framesubtitle{Student consultation}
  \archivo{set_theory_01_sample_space.py}
  \code{1}{209}{218}
  \texttt{bar\_label} delegates coordinate placement. \texttt{errorbar=None}
  avoids inferential intervals that this example did not estimate.
\end{frame}

\begin{frame}{Graphics · membership as an explicit table}
  \framesubtitle{Student consultation}
  \archivo{set_theory_02_set_operations.py}
  \textbf{Function:} \texttt{build\_membership\_figure()}.
  \code{2}{137}{144}
  The table comprehension creates one explicit 0/1 value for every outcome/set
  pair.
\end{frame}

\begin{frame}{Graphics · complete the membership table}
  \framesubtitle{Student consultation}
  \archivo{set_theory_02_set_operations.py}
  \code{2}{145}{149}
  The mathematics is complete before plotting begins.
\end{frame}

\begin{frame}{Graphics · render membership, do not infer it}
  \framesubtitle{Student consultation}
  \archivo{set_theory_02_set_operations.py}
  \code{2}{150}{155}
  Each row is a possible outcome and each column is a set. The code converts
  membership to 0/1 before Seaborn colors the matrix.
\end{frame}

\begin{frame}{Graphics · annotate the membership matrix}
  \framesubtitle{Student consultation}
  \archivo{set_theory_02_set_operations.py}
  \code{2}{156}{160}
  The figure is a membership table, not a frequency table.
\end{frame}

\begin{frame}{Graphics · annotate the membership matrix}
  \framesubtitle{Student consultation}
  \archivo{set_theory_02_set_operations.py}
  \code{2}{161}{174}
  The graphic never counts rows; it displays membership that was already computed.
\end{frame}

\begin{frame}{Graphics · reshape for a grouped comparison}
  \framesubtitle{Student consultation}
  \archivo{set_theory_03_disjoint_events.py}
  \textbf{Function:} \texttt{build\_pair\_figure()}.
  \code{3}{162}{170}
  The input table already contains the exact intersection and union sizes.
  \texttt{melt} only changes its display shape.
\end{frame}

\begin{frame}{Graphics · show two independent tests}
  \framesubtitle{Student consultation}
  \archivo{set_theory_03_disjoint_events.py}
  \code{3}{171}{176}
  The reshaped table has one row per pair and measure.
\end{frame}

\begin{frame}{Graphics · show two independent tests}
  \framesubtitle{Student consultation}
  \archivo{set_theory_03_disjoint_events.py}
  \code{3}{177}{191}
  \texttt{hue} groups the two measures; the dashed line uses the received $|S|$,
  not a graphic constant.
\end{frame}

\begin{frame}{Graphics · label the exact pair sizes}
  \framesubtitle{Student consultation}
  \archivo{set_theory_03_disjoint_events.py}
  \code{3}{192}{201}
  Labels keep zero intersection visible.
\end{frame}

\begin{frame}{Graphics · label the exact pair sizes}
  \framesubtitle{Student consultation}
  \archivo{set_theory_03_disjoint_events.py}
  \code{3}{202}{211}
  Labels keep zero intersection visible and make the two independent tests
  readable without guessing from color.
\end{frame}

\begin{frame}{Resource ownership: build, save, close}
  \framesubtitle{Student consultation}
  \archivo{set_theory_01_sample_space.py}
  \textbf{Function:} \texttt{save\_figure(figure, path)}.
  \code{1}{149}{156}
  Calculation functions return a Figure. The caller decides whether to display
  it in a notebook or save it as evidence.
\end{frame}

\begin{frame}{Resource ownership · close even on failure}
  \framesubtitle{Student consultation}
  \archivo{set_theory_01_sample_space.py}
  \code{1}{157}{162}
  \texttt{finally} closes the figure even if writing fails, so repeated
  executions do not leak resources.
\end{frame}

\begin{frame}{Validation before interpretation}
  \framesubtitle{Student consultation}
  \begin{tabular}{p{.48\linewidth}p{.45\linewidth}}
    \toprule Check & What it detects \\ \midrule
    Eight tuples and eight distinct tuples & Duplicates or omissions in $S$.\\
    Events are subsets of $S$ & Results outside the universe.\\
    Inclusion--exclusion equality & Bad filters or double counting.\\
    $A\cap A^c=\varnothing$, $A\cup A^c=S$ & Incorrect complement.\\
    Frequencies sum to 480 & Missing or duplicated rows.\\
    Seed 42 reconstructs the same rows & Accidental generator changes.\\ \bottomrule
  \end{tabular}
  \par\medskip
  A failed check requires reviewing the model and code before interpreting a figure.
\end{frame}





\end{document}
